\subsection{单连通空间}

定义 3. —— 若拓扑空间的一切覆叠都可平凡化，则称它是单连通的。

单连通空间是连通的。事实上，若空间 X 是两个非空开集 U 与 V 的不交并，则 U 到 X 中的典范单射是一个不可平凡化的覆叠。

空空间是单连通的。每个缩为一点的拓扑空间都是单连通的。

注. —— 1) 设 B 为单连通拓扑空间，(E, p) 为 B 的覆叠。若空间 E 连通且非空，则 p 是同胚。例如，设 G 为连通拓扑群，H 为 G 的离散子群，于是典范映射 p: G → G/H 使 G 成为 G/H 的覆叠（I, 第 100 页, 定理 1 的推论 2）。若空间 G/H 单连通，则 p 是同胚且 H 是平凡子群。

\begin{enumerate}
\def\labelenumi{\arabic{enumi})}
\setcounter{enumi}{1}
\tightlist
\item
  设 B 为每个点都有单连通邻域的拓扑空间；则 B 的覆叠的覆叠是 B 的覆叠。事实上，考察 B 的覆叠 (E, p) 与 E 的覆叠 (F, q)。我们来证明 (F, \(p \circ q\)) 是 B 的覆叠。由于问题是 B 上的局部问题，可设空间 B 单连通，从而 E 是 B 的可平凡化覆叠。设 V 为 E 的一个连通分支。它既开又闭，且 \(p \mid V: V \to B\) 是同胚（I, 第 69 页, 命题 1）；因此空间 V 单连通。\(q^{-1}(V)\) 的每个连通分支 W 在 \(q^{-1}(V)\) 中从而在 F 中既开又闭，且映射 q 诱导出 W 到 V 上的同胚。于是映射 \(p \circ q\) 使 F 成为 B 的可平凡化覆叠（同前引）。
\end{enumerate}

命题 3. —— 设 B 为拓扑空间。设 (E, p) 为 B 的覆叠，y 为 E 的一点。设 X 为单连通拓扑空间，f: X → B 为连续映射，x 为 X 的一点且 f(x) = p(y)。存在映射 f 的唯一的连续提升 g: X → E 使得 g(x) = y。

所求的映射 \(g\) 与连续截面 \(s\) 成一一对应（I, 第 9 页, 命题 3），这些截面是覆叠 \(\mathrm{pr}_1: \mathrm{X} \times_{\mathrm{B}} \mathrm{E} \to \mathrm{X}\) 的截面，且满足 \(s(x) = (x, y)\) 。这样的截面存在，因为空间 X 单连通；它是唯一的，因为空间 X 连通（I, 第 34 页, 命题 11 的推论 1）。

例 1. —— 设 X 为单连通拓扑空间，f 为从 X 到 \(C^{*}\) 中的连续函数。回忆（I, 第 101 页, 例 6）映射 \(z \mapsto e^{z}\) 使 C 成为 \(C^{*}\) 的覆叠。由命题 3，存在连续函数 \(h: X \to C\) 使得 \(f(x) = e^{h(x)}\) 对一切 \(x \in X\) 成立。若 \(h': X \to C\) 是另一个连续函数，满足 \(f(x) = e^{h'(x)}\) （对一切 \(x \in X\) ），则存在整数 \(q \in Z\) 使得 \(h' = h + 2\pi iq\)。

类似地，设 \(n\) 为整数 \(>0\) ；映射 \(z \mapsto z^{n}\) 使 \(C^{*}\) 成为自身的覆叠（I, 第 101 页, 例 5）。因此存在连续函数 \(k\) （从 X 到 \(C^{*}\) 中）使得 \(k(x)^{n} = f(x)\) 对一切 \(x \in X\) 成立，例如函数 \(k(x) = e^{h(x)/n}\)。若 \(k' : X \to C^{*}\) 是另一个连续函数，满足 \(k'(x)^{n} = f(x)\) （对一切 \(x \in X\) ），则存在一个 \(n\) 次单位根 \(\mu \in C\) 使得 \(k' = \mu k\)。

推论. —— 设 B 为拓扑空间，b 为 B 的一点。设 E 为 B 的单连通覆叠。对 \(E_{b}\) 的每个点 x，带基点空间 (E, x) 都是 (B, b) 的万有覆叠。

命题 4. —— 两个单连通空间（其中一个是局部连通的）的乘积是单连通空间。

设 X 与 Y 为单连通空间，且设 Y 局部连通。空间 \(X \times Y\) 是连通的，因为 X 与 Y 连通（I, 第 124 页）。设 \((Z, f)\) 为 \(X \times Y\) 的非空覆叠；我们来证明它可平凡化。由 I, 第 70 页推论 2，只需证明对 Z 的每个点 \(z_{0}\) ，存在 f 的连续截面 s 使得 \(s(f(z_{0})) = z_{0}\)。于是设 \(z_{0}\) 为 Z 的一点。置 \((x_{0}, y_{0}) = f(z_{0})\)。X × Y 的子空间 \(X \times \{y_{0}\}\) 同胚于 X。因此，由 Z 限制到 \(X \times \{y_{0}\}\) 上所得的覆叠可平凡化，且有满足 \(\sigma(x_{0}, y_{0}) = z_{0}\) 的连续截面 σ。类似地，对 X 的每个点 x，存在连续截面 \(\tau_{x}\) （f 在 \(\{x\} \times Y\) 上的截面）使得 \(\tau_{x}(x, y_{0}) = \sigma(x, y_{0})\)。对 \((x, y) \in X \times Y\) ，置 \(s(x, y) = \tau_{x}(x, y)\)。映射 s 是 f 的截面，且 \(s(x_{0}, y_{0}) = z_{0}\)。由于空间 Y 连通且局部连通，由 I, 第 35 页定理 1，映射 s 是连续的。

命题 5. —— 连通且局部连通、并且任意两个连通开子集的交都连通的拓扑空间是单连通空间。

设 B 为这样的拓扑空间。设 \((\mathrm{E}, p)\) 为 B 的覆叠，x 为 E 的一点，并记 \(b = p(x)\)。需要证明存在 p 的连续截面 s 使得 \(s(b) = x\) （推论 2, I, 第 70 页）。设 S 为偶 \((\mathrm{U}, s_{\mathrm{U}})\) 组成的集合，其中 U 是 B 的包含 b 的连通开子集，\(s_{U}\) 是 \(p_{\mathrm{U}}: p^{-1}(\mathrm{U}) \to \mathrm{U}\) 的满足 \(s_{\mathrm{U}}(b) = x\) 的连续截面。集合 S 非空（I, 第 34 页, 命题 10）。设 \((\mathrm{U}, s_{\mathrm{U}})\) 与 \((\mathrm{V}, s_{\mathrm{V}})\) 为 S 的元素。此时 \(s_{U} \mid U \cap V\) 与 \(s_{V} \mid U \cap V\) 是 \(p_{U \cap V}\) 的在 b 处取值 x 的连续截面。依假设，\(U \cap V\) 连通；因此 \(s_{U} \mid U \cap V = s_{V} \mid U \cap V\) （I, 第 34 页, 命题 11 的推论 1）。设 A 为当 \((\mathrm{U}, s_{\mathrm{U}})\) 取遍 S 时诸开集 U 的并，并设 s: \(A \to E\) 为满足 \(s \mid U = s_{U}\) （对每个偶 \((\mathrm{U}, s_{\mathrm{U}}) \in S\) ）的唯一映射。集合 A 是开且连通的（TG, I, 第 81 页, 命题 2），它包含 b，且 s 是 \(p_{A}\) 的满足 \(s(b) = x\) 的连续截面。现在只需证明集合 A 是闭的，即可得出它等于 B。

设 \(a\) 为 B 中附着于 A 的点，V 为 \(a\) 的连通开邻域，使得覆叠 E 在 V 上可平凡化。存在一点 \(c\) （属于 \(\mathrm{A} \cap \mathrm{V}\) ）以及连续截面 \(s_{\mathrm{V}}\) （\(p_{\mathrm{V}}\) 的截面），满足 \(s_{\mathrm{V}}(c) = s(c)\) 。设 \(\mathrm{A}'\) 为开集 \(\mathrm{A} \cup \mathrm{V}\)。由于 \(\mathrm{A} \cap \mathrm{V}\) 连通，存在连续截面 \(s'\) （\(p_{\mathrm{A}'}\) 的截面），它同时延拓 \(s\) 与 \(s_{\mathrm{V}}\) （I, 第 35 页, 命题 11 的推论 3）；偶 \((\mathrm{A}', s')\) 属于 \(\mathcal{S}\) ，故 \(\mathrm{A}'\) 包含于 A。于是，\(a\) 属于 A，且 A 是闭的。

推论. —— 实直线 R 的每个区间都是单连通的。

事实上，R 的连通子空间正是区间（TG, IV, 第 8 页, 定理 4），且两个区间的交是区间。

例 2. —— n 维数值空间 \(R^{n}\) （TG, VI, 第 1 页）是单连通的。事实上，它是单连通且局部连通的空间的乘积（I, 第 126 页, 命题 4 及上文命题 5 的推论）。对 \(R^{n}\) 的每个开或闭的平行六面体结论同样成立。平行体与 \(R^{n}\) 中的开或闭欧氏球都是单连通的，因为它们同胚于方块（TG, VI, 第 10 页, 命题 2 及 I, 第 124 页, 例）。

命题 6. —— 设 X 为拓扑空间。设 U₁ 与 U₂ 为 X 的开（相应地，闭）子空间，使得 X = U₁ ∪ U₂。假设 U₁ 与 U₂ 都单连通，且其交 U₁ ∩ U₂ 连通非空。则空间 X 单连通。

设 (E, p) 为 X 的覆叠，y 为 E 的一点。只需证明存在 p 的连续截面 s 使得 \(s(p(y)) = y\) （I, 第 70 页, 命题 1 的推论 2）。例如设 \(p(y)\) 属于 \(U_{1}\)。于是存在连续截面 \(s_{1}: U_{1} \to E\) （\(p_{U_{1}}\) 的截面）使得 \(s_{1}(p(y)) = y\)。设 x 为 \(U_{1} \cap U_{2}\) 的一点；存在连续截面 \(s_{2}\) （\(p_{U_{2}}\) 的截面）使得 \(s_{2}(x) = s_{1}(x)\)。由 I, 第 34 页命题 11 的推论 3，存在同时延拓 \(s_{1}\) 与 \(s_{2}\) 的 p 的连续截面 s。

例. —— 3) 对 \(n \geqslant 2\)，球面 \(\mathbf{S}_{n}\) （TG, VI, 第 10 页）是单连通的。事实上，球面 \(\mathbf{S}_{n}\) 是两个同胚于 \(\mathbf{B}_{n}\) 的闭半球之并，这两个闭半球的交同胚于

\(S_{n-1}\) （参见 TG, VI, 第 12 页）。对 \(n \geqslant 2\)，球面 \(S_{n-1}\) 连通，断言得证。

另一方面，圆 \(S_{1}\) 不是单连通的。事实上，连续映射 \(p: R \to S_{1}\) 由 \(p(\theta) = (\cos \theta, \sin \theta)\) 所定义，它使 R 成为 \(S_{1}\) 的无限次覆叠，它是连通的，因而不可平凡化。

\begin{enumerate}
\def\labelenumi{\arabic{enumi})}
\setcounter{enumi}{3}
\tightlist
\item
  设 E 为 R 上有限维 n 维向量空间，F 为 E 的余维数为 \(p \geqslant 3\) 的仿射子空间。集合 \(R^{p}\setminus\{0\}\) 同胚于 \(R \times S_{p-1}\) （TG, VI, 第 10 页, 推论 2），故它是单连通的，因为 R 与 \(S_{p-1}\) 都局部连通且单连通（I, 第 126 页, 命题 4）。集合 \(E\setminus F\) 同胚于 \(R^{n-p} \times (R^{p}\setminus\{0\})\) ，因而是单连通的（同前引）。
\end{enumerate}

命题 7. —— 设 X 为拓扑空间。设 \(U_{1}\) 与 \(U_{2}\) 为 X 的连通开（相应地，闭）子空间，使得 \(X = U_{1} \cup U_{2}\)。若空间 X 单连通，则 \(U_{1} \cap U_{2}\) 连通。

置 \(U = U_{1} \cap U_{2}\) ，并反设空间 U 不连通。我们将构造 X 的一个无限次数的连通覆叠；这样的覆叠不可平凡化。

依假设，集合 U 是两个不交非空、且在 X 中为开（相应地，闭）的集合 A 与 B 的并。对 \(i \in \{1, 2\}\) ，设 \(Y_i\) 为 \(U_i\)-空间 \((\mathrm{U}_i \times \mathbf{Z}, \mathrm{pr}_1)\) ，\(Z_i\) 为子空间 \(U \times \mathbf{Z}\) （\(Y_i\) 的子空间）。映射 \(h: Z_1 \to Z_2\) 由

\[
h (x, n) = \left\{ \begin{array}{l l} (x, n) & \text {if } x\in A \text { and } n\in \mathbf {Z} \\ (x, n + 1) & \text {if } x\in B \text { and } n\in \mathbf {Z} \end{array} \right.
\]

所定义，它是同胚。设 Y 为把 \(Y_{1}\) 与 \(Y_{2}\) 沿 \(Z_{1}\) 与 \(Z_{2}\) 借助同胚 h 粘合起来所得的空间（TG, I, 第 17 页）。

对 \(i \in \{1,2\}\) ，典范映射 \(g_i\) （从 \(Y_i\) 到 Y 中）是 \(Y_i\) 到 Y 的一个开（相应地，闭）子集上的同胚（同前引, 命题 9）。对每个整数 \(n \in \mathbf{Z}\) ，集合 \(g_i(U_i \times \{n\})\) （\(i \in \{1,2\}\)）都是连通的；由于 A 与 B 非空，\(g_1(U_1 \times \{n\})\) 与 \(g_2(U_2 \times \{n\})\) 以及 \(g_2(U_2 \times \{n+1\})\) 相交。由此得出空间 Y 连通（TG, I, 第 81 页, 推论）。

对 \(x \in \mathrm{U}\) 与 \(n \in \mathbf{Z}\) ，有 \((\mathrm{pr}_1 \circ h)(x, n) = x = \mathrm{pr}_1(x, n)\)。因此存在唯一的连续映射 \(p \colon \mathrm{Y} \to \mathrm{X}\) 使得 \(p \circ g_i = \mathrm{pr}_1\) （对 \(i \in \{1, 2\}\)）。我们来证明 X-空间 \((\mathrm{Y}, p)\) 是一个覆叠。

按构造，映射 p 的纤维同胚于离散空间 Z。

对 \(i \in \{1,2\}\) ，映射 \(g_i\) 通过过渡到子空间，定义出 \(U_i\)-空间之间从 \(\mathrm{U}_i \times \mathbf{Z}\) 到 \(p^{-1}(\mathrm{U}_i)\) 上的一个同构。

按空间 Y 的定义，存在唯一的从 \((\mathrm{X}\setminus\mathrm{A})\times\mathbf{Z}\) 到 Y 中的映射 k，使得 \(k(x,n)=g_{1}(x,n)\) （当 \(x\in(\mathrm{U}_{1}\setminus\mathrm{A})\times\mathbf{Z}\) ），\(k(x,n)=g_{2}(x,n-1)\) （当 \(x\in(\mathrm{U}_{2}\setminus\mathrm{A})\times\mathbf{Z}\) ）；它是 \((\mathrm{X}\setminus\mathrm{A})\)-空间之间从 \((\mathrm{X}\setminus\mathrm{A})\times\mathbf{Z}\) 到 \(p^{-1}(\mathrm{X}\setminus\mathrm{A})\) 上的同构。同样，存在唯一的映射 \(k'\) （从 \((\mathrm{X}\setminus\mathrm{B})\times\mathbf{Z}\) 到 Y 中），它与 \(g_{1}\) 在 \((\mathrm{U}_{1}\setminus\mathrm{B})\times\mathbf{Z}\) 上相合，与 \(g_{2}\) 在 \((\mathrm{U}_{2}\setminus\mathrm{B})\times\mathbf{Z}\) 上相合；它是 \((\mathrm{X}\setminus\mathrm{B})\)-空间之间从 \((\mathrm{X}\setminus\mathrm{B})\times\mathbf{Z}\) 到 \(p^{-1}(\mathrm{X}\setminus\mathrm{B})\) 上的同构。

这就证明 X-空间 \((Y,p)\) 在部分 \(U_{1}\)、\(U_{2}\)、\(X\setminus A\) 与 \(X\setminus B\) 上都可平凡化。我们有 \(U_{1} \cup U_{2} = X\) ；若 \(U_{1}\) 与 \(U_{2}\) 在 X 中是开的，这就证明 \((Y,p)\) 是 X 的覆叠。当 \(U_{1}\) 与 \(U_{2}\) 在 X 中是闭时结论同样成立，因为此时 \(X\setminus A\) 与 \(X\setminus B\) 是 X 的开子集且其并为 X。命题至此得证。

推论. —— 设 X 为单连通拓扑空间，A 为 X 的连通子集。若 A 的余集连通，则其边界也连通。

设 \(X_{1}\) 与 \(X_{2}\) 分别为 A 与 \(X\setminus A\) 的闭包。集合 \(X_{1}\) 与 \(X_{2}\) 是闭且连通的（TG, I, 第 81 页, 命题 1）；我们有 \(X_{1} \cup X_{2} = X\) ，且其交 \(X_{1} \cap X_{2}\) 等于 A 的边界。于是只需应用命题 7。
% semantic-fragments: legacy-input-seam
\relax{}
